\documentclass[a4paper,11pt]{article}

% ── Codifica e lingua ────────────────────────────────────────────────
\usepackage[utf8]{inputenc}
\usepackage[T1]{fontenc}
\usepackage[italian]{babel}

% ── Font e tipografia ────────────────────────────────────────────────
\usepackage{lmodern}
\usepackage{microtype}

% ── Layout pagina ────────────────────────────────────────────────────
\usepackage[
  top=28mm, bottom=26mm,
  left=26mm, right=26mm,
  headheight=14pt
]{geometry}

% ── Intestazioni e piè di pagina ─────────────────────────────────────
\usepackage{fancyhdr}
\pagestyle{fancy}
\fancyhf{}
\fancyhead[L]{\small\textsc{StesEye} -- Relazione Tecnica}
\fancyhead[R]{\small\thepage}
\fancyfoot[C]{\footnotesize Documento interno}
\renewcommand{\headrulewidth}{0.4pt}
\renewcommand{\footrulewidth}{0.2pt}

% ── Pacchetti vari ───────────────────────────────────────────────────
\usepackage{graphicx}
\usepackage{xcolor}
\usepackage{booktabs}
\usepackage{tabularx}
\usepackage{longtable}
\usepackage{enumitem}
\usepackage{listings}
\usepackage{caption}
\usepackage{float}
\usepackage{parskip}
\usepackage{tikz}
\usepackage[hidelinks]{hyperref}

% ── Colori ───────────────────────────────────────────────────────────
\definecolor{accent}{HTML}{388bfd}
\definecolor{accentdark}{HTML}{1a4a8a}
\definecolor{darkbg}{HTML}{0d1117}
\definecolor{cardbg}{HTML}{161b22}
\definecolor{codebg}{HTML}{f5f5f5}
\definecolor{codeframe}{HTML}{d0d0d0}
\definecolor{greenst}{HTML}{3fb950}
\definecolor{greendim}{HTML}{1a3a2a}
\definecolor{reddim}{HTML}{f85149}
\definecolor{yellowst}{HTML}{d29922}
\definecolor{lightgray}{HTML}{e6edf3}
\definecolor{mutedgray}{HTML}{7d8590}

% ── Stile listings ───────────────────────────────────────────────────
\lstdefinestyle{pythonstyle}{
  language=Python,
  backgroundcolor=\color{codebg},
  frame=single,
  rulecolor=\color{codeframe},
  basicstyle=\ttfamily\small,
  keywordstyle=\bfseries\color{accent},
  commentstyle=\itshape\color{gray},
  stringstyle=\color{teal},
  showstringspaces=false,
  breaklines=true,
  tabsize=4,
  xleftmargin=4mm,
  framexleftmargin=3mm,
  numbers=left,
  numberstyle=\tiny\color{gray},
  numbersep=6pt,
  captionpos=b
}
\lstset{style=pythonstyle}

% ── Sezione con riga colorata ────────────────────────────────────────
\usepackage{titlesec}
\titleformat{\section}
  {\large\bfseries}{\thesection.}{0.5em}{}[\vspace{2pt}\textcolor{accent}{\rule{\textwidth}{0.6pt}}]
\titleformat{\subsection}
  {\normalsize\bfseries}{\thesubsection}{0.5em}{}

% ── Comando helper: box colorato in-linea ────────────────────────────
\newcommand{\infobox}[2]{%
  \begin{tikzpicture}[baseline=(text.base)]
    \node[fill=#1, rounded corners=3pt, inner sep=5pt, text=white,
          font=\small\bfseries] (text) {#2};
  \end{tikzpicture}%
}

% ═════════════════════════════════════════════════════════════════════
\begin{document}

% ═════════════════════════════════════════════════════════════════════
%  FRONTESPIZIO CUSTOM
% ═════════════════════════════════════════════════════════════════════
\thispagestyle{empty}
\begin{tikzpicture}[remember picture, overlay]

  % --- Sfondo scuro in alto (60% della pagina) ---
  \fill[darkbg]
    (current page.north west) rectangle
    ([yshift=-0.58\paperheight]current page.north east);

  % --- Banda accent diagonale ---
  \fill[accent, opacity=0.15]
    ([yshift=-0.42\paperheight, xshift=-10mm]current page.north west)
    -- ([yshift=-0.52\paperheight, xshift=-10mm]current page.north west)
    -- ([yshift=-0.48\paperheight]current page.north east)
    -- ([yshift=-0.38\paperheight]current page.north east)
    -- cycle;

  % --- Cerchio decorativo grande ---
  \fill[accent, opacity=0.08]
    ([xshift=40mm, yshift=-70mm]current page.north east) circle (80mm);

  % --- Cerchio decorativo piccolo ---
  \fill[accent, opacity=0.12]
    ([xshift=-30mm, yshift=-160mm]current page.north east) circle (25mm);

  % --- Linea accent orizzontale ---
  \draw[accent, line width=2pt]
    ([xshift=26mm, yshift=-92mm]current page.north west)
    -- ([xshift=-26mm, yshift=-92mm]current page.north east);

  % --- Titolo principale ---
  \node[anchor=west, text=lightgray, font=\fontsize{36}{42}\bfseries\selectfont]
    at ([xshift=30mm, yshift=-50mm]current page.north west)
    {StesEye};

  % --- Pallino accent (come nel logo dell'app) ---
  \fill[accent]
    ([xshift=136mm, yshift=-43mm]current page.north west) circle (4pt);

  % --- Sottotitolo ---
  \node[anchor=west, text=mutedgray, font=\fontsize{13}{18}\selectfont]
    at ([xshift=30mm, yshift=-64mm]current page.north west)
    {Network Scanner \& Device Analyzer};

  % --- Label "Relazione Tecnica" ---
  \node[anchor=west, text=accent, font=\fontsize{16}{22}\bfseries\selectfont]
    at ([xshift=30mm, yshift=-80mm]current page.north west)
    {Relazione Tecnica};

  % --- Box informativi sotto la linea ---
  % Box 1: Linguaggio
  \node[anchor=north west, fill=cardbg, rounded corners=6pt,
        minimum width=72mm, minimum height=28mm,
        inner sep=0pt, outer sep=0pt]
    (box1) at ([xshift=26mm, yshift=-105mm]current page.north west) {};
  \node[anchor=north west, text=mutedgray, font=\footnotesize\bfseries]
    at ([xshift=4mm, yshift=-4mm]box1.north west) {LINGUAGGIO};
  \node[anchor=north west, text=lightgray, font=\large\bfseries]
    at ([xshift=4mm, yshift=-12mm]box1.north west) {Python 3.10+};
  \node[anchor=north west, text=mutedgray, font=\footnotesize]
    at ([xshift=4mm, yshift=-21mm]box1.north west) {$\sim$2\,900 righe in 13 moduli};

  % Box 2: Piattaforma
  \node[anchor=north west, fill=cardbg, rounded corners=6pt,
        minimum width=72mm, minimum height=28mm,
        inner sep=0pt, outer sep=0pt]
    (box2) at ([xshift=104mm, yshift=-105mm]current page.north west) {};
  \node[anchor=north west, text=mutedgray, font=\footnotesize\bfseries]
    at ([xshift=4mm, yshift=-4mm]box2.north west) {PIATTAFORMA};
  \node[anchor=north west, text=lightgray, font=\large\bfseries]
    at ([xshift=4mm, yshift=-12mm]box2.north west) {Windows 11};
  \node[anchor=north west, text=mutedgray, font=\footnotesize]
    at ([xshift=4mm, yshift=-21mm]box2.north west) {Privilegi di amministratore};

  % Box 3: GUI
  \node[anchor=north west, fill=cardbg, rounded corners=6pt,
        minimum width=72mm, minimum height=28mm,
        inner sep=0pt, outer sep=0pt]
    (box3) at ([xshift=26mm, yshift=-140mm]current page.north west) {};
  \node[anchor=north west, text=mutedgray, font=\footnotesize\bfseries]
    at ([xshift=4mm, yshift=-4mm]box3.north west) {INTERFACCIA};
  \node[anchor=north west, text=lightgray, font=\large\bfseries]
    at ([xshift=4mm, yshift=-12mm]box3.north west) {CustomTkinter};
  \node[anchor=north west, text=mutedgray, font=\footnotesize]
    at ([xshift=4mm, yshift=-21mm]box3.north west) {Dark theme nativo};

  % Box 4: Versione
  \node[anchor=north west, fill=cardbg, rounded corners=6pt,
        minimum width=72mm, minimum height=28mm,
        inner sep=0pt, outer sep=0pt]
    (box4) at ([xshift=104mm, yshift=-140mm]current page.north west) {};
  \node[anchor=north west, text=mutedgray, font=\footnotesize\bfseries]
    at ([xshift=4mm, yshift=-4mm]box4.north west) {VERSIONE};
  \node[anchor=north west, text=lightgray, font=\large\bfseries]
    at ([xshift=4mm, yshift=-12mm]box4.north west) {1.0};
  \node[anchor=north west, text=mutedgray, font=\footnotesize]
    at ([xshift=4mm, yshift=-21mm]box4.north west) {Maggio 2026};

  % --- Tre feature pills ---
  \node[fill=greendim, rounded corners=4pt, inner xsep=8pt, inner ysep=4pt,
        text=greenst, font=\small\bfseries, anchor=west]
    (f1) at ([xshift=26mm, yshift=-182mm]current page.north west)
    {Discovery};
  \node[fill=accentdark!60, rounded corners=4pt, inner xsep=8pt, inner ysep=4pt,
        text=accent, font=\small\bfseries, anchor=west]
    at ([xshift=6mm]f1.east)
    {Analysis};
  \node[fill=greendim, rounded corners=4pt, inner xsep=8pt, inner ysep=4pt,
        text=greenst, font=\small\bfseries, anchor=west]
    at ([xshift=58mm]f1.east)
    {File Transfer};

  % --- Descrizione breve ---
  \node[anchor=north west, text width=140mm, text=black!80,
        font=\normalsize, align=left]
    at ([xshift=26mm, yshift=-0.65\paperheight]current page.north west)
    {Applicazione desktop per la scansione automatica delle reti locali.
     Individua i dispositivi connessi, ne riconosce il tipo e il sistema
     operativo, e permette di inviare file via SMB verso le macchine
     Windows della rete.};

  % --- Info in basso ---
  \node[anchor=south west, text=black!50, font=\footnotesize]
    at ([xshift=26mm, yshift=22mm]current page.south west)
    {Dipendenze: tutte open-source (MIT / BSD)};

  \node[anchor=south east, text=black!50, font=\footnotesize]
    at ([xshift=-26mm, yshift=22mm]current page.south east)
    {Maggio 2026};

\end{tikzpicture}

\newpage

% ═════════════════════════════════════════════════════════════════════
%  INDICE
% ═════════════════════════════════════════════════════════════════════
\tableofcontents
\newpage

% ═════════════════════════════════════════════════════════════════════
\section{Introduzione}

StesEye è un programma desktop che serve a scoprire e tenere
sotto controllo tutti i dispositivi collegati a una rete locale
(LAN).  L'idea di base è semplice: si lancia l'app, si preme
``Scan'' e in pochi secondi compare la lista di tutto ciò che è
connesso alla rete -- PC, telefoni, stampanti, telecamere, router,
NAS e così via.

È pensato per situazioni in cui si vuole capire velocemente
``chi c'è in rete'' senza installare software pesanti o
configurare strumenti da riga di comando.  Funziona bene in reti
domestiche, laboratori scolastici, piccoli uffici.

Il software gira su Windows~11 e ha bisogno dei privilegi di
amministratore perché deve mandare pacchetti ARP sulla rete
(oppure usare il comando \texttt{ping}), operazioni che
Windows riserva agli utenti con permessi elevati.

Le tre funzionalità principali sono:

\begin{enumerate}[nosep]
  \item \textbf{Discovery} -- trova automaticamente tutti i
        dispositivi connessi alla sottorete.
  \item \textbf{Analysis} -- per ogni dispositivo cerca di capire
        che cos'è (router, PC, stampante\ldots), che sistema
        operativo usa e quali servizi ha attivi.
  \item \textbf{File Transfer} -- permette di inviare file verso
        PC Windows attraverso le cartelle condivise (SMB).
\end{enumerate}


% ═════════════════════════════════════════════════════════════════════
\section{Requisiti e dipendenze}

\subsection{Cosa serve per farlo funzionare}
\begin{itemize}[nosep]
  \item Windows 11 (oppure Windows 10 aggiornato).
  \item Python 3.10 o più recente.
  \item Un account con privilegi di amministratore.
  \item Per la scansione ARP veloce: il driver Npcap installato
        (si scarica gratis dal sito ufficiale).
\end{itemize}

\subsection{Librerie Python}
L'unica libreria obbligatoria è \texttt{customtkinter} (il
framework per la grafica).  Le altre due sono facoltative: se
ci sono il programma funziona meglio, ma anche senza va
benissimo.

\begin{table}[H]
\centering
\small
\begin{tabularx}{0.92\textwidth}{llX}
  \toprule
  \textbf{Pacchetto}       & \textbf{Obbligatorio?} & \textbf{A cosa serve} \\
  \midrule
  \texttt{customtkinter}   & sì & interfaccia grafica con tema scuro \\
  \texttt{scapy}           & no & manda e riceve pacchetti ARP grezzi
                                  (scan veloce) \\
  \texttt{mac-vendor-lookup}& no & dato un indirizzo MAC, dice il
                                   produttore (es.\ ``Apple'', ``TP-Link'') \\
  \bottomrule
\end{tabularx}
\end{table}

Quando Scapy non è presente, il programma passa in automatico
al metodo alternativo: manda un \texttt{ping} a tutti e poi
legge la tabella ARP di sistema con \texttt{arp~-a}.  È un po'
più lento, ma funziona senza bisogno di driver aggiuntivi.


% ═════════════════════════════════════════════════════════════════════
\section{Architettura del progetto}

Il codice è organizzato come un package Python (la cartella
\texttt{steseye/}).  Ogni file ha un compito preciso: questo
rende più facile leggere il codice, trovare i bug e lavorarci
sopra.

\begin{table}[H]
\centering
\small
\begin{tabularx}{\textwidth}{llX}
  \toprule
  \textbf{File} & \textbf{Righe} & \textbf{Cosa fa} \\
  \midrule
  \texttt{\_\_init\_\_.py}   &    1 & Dice a Python che questa cartella
                                      è un package \\
  \texttt{\_\_main\_\_.py}   &   39 & Punto di partenza: controlla i
                                      permessi e avvia l'app \\
  \texttt{colors.py}         &   31 & Tutti i colori dell'interfaccia
                                      in un unico posto \\
  \texttt{device\_types.py}  &   92 & Tabelle per riconoscere i tipi di
                                      dispositivo \\
  \texttt{device.py}         &   33 & La ``scheda'' di un dispositivo
                                      (IP, MAC, tipo, ecc.) \\
  \texttt{analyzer.py}       &  230 & Il cervello che classifica i
                                      dispositivi + port scanner \\
  \texttt{file\_transfer.py} &  228 & Tutta la logica per mandare file
                                      via cartelle condivise \\
  \texttt{scan\_engine.py}   &  332 & Il motore che scopre i dispositivi
                                      in rete (ARP o ping) \\
  \texttt{widgets.py}        &  222 & Componenti grafici riutilizzabili
                                      (logo, pallino, riga dispositivo) \\
  \texttt{detail\_panel.py}  &  202 & Il pannello laterale con i
                                      dettagli del device selezionato \\
  \texttt{dialogs.py}        &  943 & Le finestre pop-up (risultato
                                      invio, browser cartelle, ecc.) \\
  \texttt{app.py}            &  531 & La finestra principale con tutta
                                      la logica di interfaccia \\
  \texttt{run.py}            &    7 & Un piccolo launcher di comodo \\
  \bottomrule
\end{tabularx}
\caption{I moduli del package \texttt{steseye}.}
\end{table}

\subsection{Come si collegano tra loro i moduli}

I file si importano a vicenda secondo uno schema a livelli:
i moduli ``bassi'' (colori, tipi, dati) non importano nessun
altro, quelli ``alti'' (app, main) importano tutti gli altri.
Non ci sono dipendenze circolari.

\begin{center}
\small
\begin{verbatim}
  colors.py
     |
  device_types.py     device.py
     |       \__________|
     |                  |
  analyzer.py      file_transfer.py
     |                  |
  scan_engine.py        |
     |                  |
  widgets.py   detail_panel.py   dialogs.py
     \____________|_____________/
                  |
               app.py
                  |
            __main__.py
\end{verbatim}
\end{center}

\subsection{Avvio del programma}

Per lanciare StesEye basta aprire un terminale e scrivere:

\begin{lstlisting}[numbers=none]
python -m steseye
\end{lstlisting}

Il file \texttt{\_\_main\_\_.py} come prima cosa controlla se
l'utente ha i permessi di amministratore.  Lo fa chiamando la
funzione \texttt{IsUserAnAdmin()} delle API di Windows.
Se non siamo admin, compare una finestra che chiede se rilanciare
il programma ``come amministratore'' (la classica finestra UAC
di Windows).


% ═════════════════════════════════════════════════════════════════════
\section{Motore di scansione}

Il motore di scansione è il pezzo più importante del programma.
Sta tutto nella classe \texttt{ScanEngine} dentro
\texttt{scan\_engine.py}.  Il suo lavoro si divide in due fasi:
prima scopre \emph{chi} c'è in rete, poi \emph{analizza} ogni
dispositivo trovato.

\subsection{Fase 1 -- Scoperta dei dispositivi (Discovery)}

All'avvio, il motore sceglie il metodo di scansione migliore
tra due opzioni:

\begin{description}[style=nextline, leftmargin=2.5cm]
  \item[Metodo ARP (Scapy)]
    Manda un singolo pacchetto ARP broadcast a tutta la sottorete
    /24 (cioè 254 indirizzi) e aspetta le risposte.
    Chi risponde è connesso.  È il metodo più veloce: ci mette
    2--4 secondi per tutta la rete.  Per sicurezza fa due passate
    consecutive, così se un dispositivo non ha risposto al primo
    tentativo c'è una seconda possibilità.

  \item[Metodo Fallback (Ping + ARP)]
    Se Scapy o Npcap non sono installati, il programma lancia
    in parallelo 254 comandi \texttt{ping} (uno per ogni IP
    possibile), usando fino a 50 thread contemporanei.  Dopo
    che i ping sono finiti, legge la tabella ARP del sistema
    (con il comando \texttt{arp~-a}) dove Windows ha salvato
    le risposte ricevute.
\end{description}

Dopo la scansione, il motore confronta i risultati con la lista
dei dispositivi che già conosce:
\begin{itemize}[nosep]
  \item se un MAC è nuovo, crea una nuova scheda dispositivo;
  \item se lo conosceva già, aggiorna l'orario di ``ultimo
        contatto'';
  \item se un dispositivo che prima c'era adesso non risponde,
        lo segna come ``Offline''.
\end{itemize}

La lista è protetta da un \texttt{threading.Lock} perché più
thread ci lavorano sopra contemporaneamente.

\subsection{Fase 2 -- Analisi dei dispositivi}

Per ogni dispositivo appena trovato (o su cui l'utente chiede
una ri-analisi), parte un thread che raccoglie quante più
informazioni possibili:

\begin{enumerate}[nosep]
  \item \textbf{Port scan} -- prova a connettersi a 21 porte
        TCP predefinite (SSH, HTTP, SMB, RDP\ldots).
        Per ogni porta usa un timeout di 0.6 secondi, con
        un massimo di 25 tentativi in parallelo.
  \item \textbf{TTL} -- manda un \texttt{ping} e guarda il
        campo Time-To-Live della risposta: valori diversi
        suggeriscono sistemi operativi diversi.
  \item \textbf{Banner HTTP} -- se la porta 80, 8080 o 443 è
        aperta, manda una richiesta \texttt{HEAD} e legge
        l'header \texttt{Server:} (es.\ ``nginx'', ``Synology'').
  \item \textbf{NetBIOS} -- se la porta 445 è aperta o il TTL
        suggerisce Windows, manda una query UDP sulla porta 137
        per scoprire il nome macchina nella rete Windows.
\end{enumerate}

Tutti questi dati vengono poi passati al classificatore
(descritto nella prossima sezione) che decide ``che cos'è''
quel dispositivo.

\subsection{Monitoraggio continuo}

La modalità \emph{Monitor} ripete la scansione a intervalli
regolari (di default ogni 30 secondi, ma si può cambiare).
Il thread di monitoraggio si può interrompere in qualsiasi
momento premendo ``Stop'': internamente usa un
\texttt{threading.Event} che viene controllato ogni mezzo
secondo.

Ad ogni ciclo, i dispositivi che non rispondono più vengono
colorati di rosso nell'interfaccia.


% ═════════════════════════════════════════════════════════════════════
\section{Sistema di classificazione}

La classe \texttt{Analyzer} (\texttt{analyzer.py}) è quella che
decide se un dispositivo è un router, un PC, una stampante,
un telefono, eccetera.  Lo fa con un \textbf{sistema a punti}:
per ogni indizio trovato assegna dei punti a un certo tipo di
dispositivo, e alla fine vince il tipo che ha accumulato più
punti.

\subsection{Gli indizi usati (euristiche)}

\begin{table}[H]
\centering
\small
\begin{tabularx}{\textwidth}{lrX}
  \toprule
  \textbf{Indizio} & \textbf{Punti} & \textbf{Come funziona} \\
  \midrule
  Produttore (OUI) & fino a 30  & Il MAC address contiene un codice del
                                   produttore.  Se il vendor è ``TP-Link''
                                   probabilmente è un router; se è ``Canon''
                                   probabilmente è una stampante. \\[3pt]
  Porte aperte      & 15--40     & Ogni porta aperta suggerisce un tipo.
                                   Porta 3389 aperta? Probabilmente è un
                                   PC Windows con Remote Desktop.
                                   Porte 9100+515? Stampante.  E così via.
                                   Alcune combinazioni danno bonus extra. \\[3pt]
  TTL               & 10--15     & TTL $\leq$ 64 di solito significa Linux o
                                   macOS; tra 65 e 128 è quasi sempre
                                   Windows; sopra 128 è un apparato di rete. \\[3pt]
  Nome host         & 20--45     & Se il nome contiene ``iphone'' è un iPhone;
                                   ``printer'' è una stampante; ``raspberrypi''
                                   è un Raspberry Pi.  Ogni keyword ha un
                                   punteggio diverso. \\[3pt]
  Default gateway   & 50         & Se l'IP del dispositivo coincide con il
                                   gateway (il ``router di uscita''), prende
                                   50 punti come Router.  È il bonus più
                                   alto perché è quasi certo. \\[3pt]
  Banner HTTP / NetBIOS & 15--35 & Se il server web risponde ``Synology'' è un NAS; ``Hikvision'' è una telecamera. Un nome NetBIOS presente indica Windows. \\
  \bottomrule
\end{tabularx}
\caption{Le sei euristiche del classificatore con i pesi associati.}
\end{table}

Il punteggio finale è limitato a 100 e diventa il valore di
\emph{confidence} (confidenza) che si vede nell'interfaccia: più
è alto, più il programma è ``sicuro'' di aver indovinato il tipo.

\subsection{Badge nella lista}

Ad ogni tipo di dispositivo corrisponde un'etichetta corta di
2--3 lettere e un colore, usati come badge nella lista
dispositivi.  Qualche esempio:

\begin{center}
\small
\infobox{accent}{RTR} Router \qquad
\infobox{greenst}{WIN} Windows PC \qquad
\infobox{orange}{PRT} Stampante \qquad
\infobox{reddim}{CAM} Telecamera \qquad
\infobox{yellowst}{NAS} NAS
\end{center}


% ═════════════════════════════════════════════════════════════════════
\section{Trasferimento file via SMB}

Il modulo \texttt{file\_transfer.py} permette di mandare file
a un altro PC Windows della rete, usando le cartelle condivise
(protocollo SMB, cioè quello che Windows usa normalmente per
condividere file in rete locale).

\subsection{Come trova le cartelle condivise}

Il programma prova tre metodi diversi, in cascata:

\begin{enumerate}[nosep]
  \item Esegue il comando \texttt{net view} e legge l'elenco
        delle share (cartelle condivise) di tipo ``Disk''.
  \item Se il primo metodo non funziona, prova con una
        query WMI via PowerShell.
  \item Come ultimo tentativo, prova ad accedere direttamente
        ad alcune share note (\texttt{C\$}, \texttt{Users},
        \texttt{Public}).
\end{enumerate}

\subsection{Come copia i file}

La copia avviene con le normali funzioni di I/O di Python:
\texttt{shutil.copy2} per file piccoli; lettura e scrittura a
blocchi da 1~MB per file grandi (così la barra di progresso si
aggiorna man mano).

Se nella cartella di destinazione esiste già un file con lo
stesso nome, il programma aggiunge un numero tra parentesi
(es.\ \texttt{relazione~(1).pdf}) invece di sovrascriverlo.

\subsection{Interfaccia di invio}

La finestra ``Send File'' guida l'utente in due passi:

\begin{enumerate}[nosep]
  \item Si scelgono i file da inviare (con il classico dialogo
        ``Sfoglia'').
  \item Si sceglie la destinazione: si può cliccare su una share
        dall'elenco, navigare le sottocartelle con un browser
        integrato, oppure digitare manualmente un percorso UNC
        (tipo \verb|\\192.168.1.50\Users\Public|).
\end{enumerate}

Alla fine compare un riepilogo con l'esito di ogni file: quali
sono stati inviati, quali hanno dato errore, e un pulsante per
aprire la cartella di destinazione.


% ═════════════════════════════════════════════════════════════════════
\section{Interfaccia grafica}

L'interfaccia è fatta con CustomTkinter, un framework che
estende Tkinter (quello standard di Python) con uno stile
moderno e il supporto al tema scuro.

I colori sono tutti definiti in un unico file
(\texttt{colors.py}): 26 costanti esadecimali che coprono
sfondi, testi, colori di accento, stati online/offline/warning.
La palette è ispirata al tema scuro di GitHub.

\subsection{Le quattro aree della finestra}

\begin{description}[style=nextline, leftmargin=2.5cm]
  \item[Top bar]
    Contiene il logo, il campo dove si inserisce il range di rete
    (es.\ \texttt{192.168.1.0/24}), il timeout, l'intervallo di
    monitoraggio e i pulsanti Scan / Monitor / Stop / Clear /
    Export.  C'è anche un checkbox ``Alerts'' per attivare i
    suoni di notifica.

  \item[Lista dispositivi]
    L'area principale: una lista scrollabile dove ogni riga è
    una ``card'' che mostra il pallino verde/rosso di stato,
    il badge del tipo, il nome, l'IP, il MAC, il vendor, il
    sistema operativo stimato, la percentuale di confidenza
    e l'orario dell'ultimo rilevamento.
    Cliccando con il tasto destro si apre un menu contestuale
    per copiare dati, ri-analizzare o inviare file.

  \item[Pannello dettaglio]
    Un pannello fisso sulla destra (largo 340~pixel) che mostra
    tutti i dati del dispositivo selezionato, organizzati in
    sezioni: Network, Identity, Services, Analysis, Timeline.
    In fondo ci sono i pulsanti ``Re-analyze'', ``Copy info'' e
    (solo per i PC Windows con la porta 445 aperta) ``Send file''.

  \item[Barra di stato]
    In basso: contatori (online / totale / analizzati), una barra
    di progresso che si anima durante la scansione e un log
    testuale con orario di ogni evento.
\end{description}

\subsection{Thread e GUI: come convivono}

Tutte le operazioni di rete (scansione, analisi, scoperta share,
trasferimento file) girano su thread separati per non bloccare
l'interfaccia.  L'aggiornamento dei componenti grafici avviene
tramite \texttt{widget.after(0, callback)}, che accoda l'esecuzione
nel loop principale di Tkinter.  Questo è necessario perché Tkinter
non è thread-safe: se si aggiornasse un widget da un thread
secondario, il programma potrebbe andare in crash.


% ═════════════════════════════════════════════════════════════════════
\section{La classe \texttt{Device}}

Ogni dispositivo trovato in rete è rappresentato da un oggetto
della classe \texttt{Device} (file \texttt{device.py}).  Questa
classe usa \texttt{\_\_slots\_\_}, una tecnica di Python che
riduce il consumo di memoria forzando la lista degli attributi
possibili.

\begin{table}[H]
\centering
\small
\begin{tabularx}{\textwidth}{llX}
  \toprule
  \textbf{Attributo}    & \textbf{Tipo}   & \textbf{Cosa contiene} \\
  \midrule
  \texttt{ip}           & str             & Indirizzo IPv4 del dispositivo \\
  \texttt{mac}          & str             & Indirizzo MAC (es.\ \texttt{aa:bb:cc:11:22:33}) \\
  \texttt{vendor}       & str             & Nome del produttore (da OUI lookup) \\
  \texttt{hostname}     & str             & Nome host (ottenuto con reverse DNS) \\
  \texttt{first\_seen}  & datetime        & Quando è stato visto la prima volta \\
  \texttt{last\_seen}   & datetime        & Quando ha risposto l'ultima volta \\
  \texttt{status}       & str             & ``Online'' oppure ``Offline'' \\
  \texttt{device\_type} & str             & Tipo classificato (es.\ ``Router'', ``Printer'') \\
  \texttt{os\_guess}    & str             & Sistema operativo stimato \\
  \texttt{open\_ports}  & list[int]       & Porte TCP risultate aperte \\
  \texttt{services}     & list[str]       & Servizi trovati (es.\ ``80/HTTP'') \\
  \texttt{ttl}          & int             & Time-To-Live dall'ultimo ping \\
  \texttt{http\_banner} & str             & Header \texttt{Server:} della risposta HTTP \\
  \texttt{netbios\_name}& str             & Nome NetBIOS (reti Windows) \\
  \texttt{confidence}   & int             & Punteggio di confidenza (da 0 a 100) \\
  \texttt{badge}        & str             & Etichetta corta per la UI (es.\ ``RTR'') \\
  \texttt{analyzed}     & bool            & \texttt{True} se l'analisi è stata completata \\
  \bottomrule
\end{tabularx}
\caption{Tutti gli attributi della classe \texttt{Device}.}
\end{table}


% ═════════════════════════════════════════════════════════════════════
\section{Esportazione dei risultati}

I dispositivi trovati si possono esportare in due formati tramite
il pulsante ``Export'':

\begin{itemize}[nosep]
  \item \textbf{JSON} -- un file con un array di oggetti, ognuno
        con tutti i campi del dispositivo (IP, MAC, vendor,
        hostname, tipo, OS, porte aperte, confidence, status,
        timestamp).  Comodo per elaborazioni successive con
        script o per importarlo in altri programmi.
  \item \textbf{CSV} -- un file tabellare con una riga per
        dispositivo e le colonne separate da virgola.
        Si apre direttamente con Excel o LibreOffice Calc.
\end{itemize}

In entrambi i casi il file è codificato in UTF-8.


% ═════════════════════════════════════════════════════════════════════
\section{Elenco delle porte scansionate}

Il port scan controlla 21 porte TCP, scelte per coprire i servizi
che si trovano più spesso in reti domestiche e d'ufficio.

\begin{table}[H]
\centering
\small
\begin{tabular}{rll}
  \toprule
  \textbf{Porta} & \textbf{Servizio}   & \textbf{Tipo suggerito} \\
  \midrule
  22    & SSH           & Server \\
  23    & Telnet        & Apparato di rete \\
  53    & DNS           & DNS Server \\
  80    & HTTP          & Web Server \\
  443   & HTTPS         & Web Server \\
  445   & SMB           & File Share / Windows \\
  515   & LPR           & Stampante \\
  548   & AFP           & Mac / NAS \\
  554   & RTSP          & Telecamera \\
  631   & IPP           & Stampante \\
  3389  & RDP           & Windows PC \\
  5000  & HTTP          & NAS (Synology) \\
  5353  & mDNS          & Dispositivo Apple \\
  5900  & VNC           & Remote Desktop \\
  7000  & AirPlay       & Apple TV \\
  8008  & HTTP alt.     & Chromecast \\
  8009  & Cast protocol & Chromecast \\
  8080  & HTTP alt.     & Server \\
  9100  & RAW Print     & Stampante \\
  32400 & Plex          & Media Server \\
  62078 & Lockdown      & iOS \\
  \bottomrule
\end{tabular}
\caption{Le 21 porte scansionate con il servizio corrispondente.}
\end{table}


% ═════════════════════════════════════════════════════════════════════
\section{Limiti e aspetti di sicurezza}

Come ogni software, StesEye ha dei limiti di cui è bene essere
consapevoli:

\begin{itemize}[nosep]
  \item Servono i \textbf{privilegi di amministratore}, quindi non
        si può usare su PC dove non si hanno questi permessi.
  \item La scansione ARP funziona solo all'interno della stessa
        sottorete (stesso dominio di broadcast L2): \textbf{non
        attraversa router o VLAN}.
  \item L'invio di file SMB usa le credenziali con cui si è
        loggati su Windows: se non si hanno i permessi sulla
        macchina remota, il trasferimento fallisce.
  \item I file vengono mandati \textbf{senza cifratura aggiuntiva}:
        la sicurezza dipende da come è configurato SMB sulla rete.
  \item Si scansionano solo 21 porte fisse.  Un dispositivo che
        usa porte non standard potrebbe non essere riconosciuto
        correttamente.
\end{itemize}


% ═════════════════════════════════════════════════════════════════════
\section{Sviluppi futuri}

Alcune idee per le prossime versioni:

\begin{enumerate}[nosep]
  \item Supporto anche per \textbf{macOS e Linux}, adattando i
        comandi di sistema (\texttt{ping}, \texttt{arp}).
  \item \textbf{Salvataggio su database} (SQLite) per tenere uno
        storico delle scansioni e confrontare le reti nel tempo.
  \item Possibilità di \textbf{scegliere le porte} da scansionare
        direttamente dalla GUI.
  \item \textbf{Notifiche push} (Telegram, Slack) quando compare
        un dispositivo sconosciuto in rete.
  \item Integrazione di \textbf{Nmap} come alternativa per il
        port scan, che è più preciso nel riconoscere i servizi.
  \item \textbf{Generazione PDF} dei report direttamente
        dall'interfaccia.
\end{enumerate}


% ═════════════════════════════════════════════════════════════════════
\section{Conclusioni}

StesEye è un tool che fa bene il suo lavoro: in pochi secondi
mostra una mappa chiara di tutti i dispositivi collegati alla rete
locale, con informazioni utili come il tipo di dispositivo, il
sistema operativo e i servizi attivi.  Il sistema a punteggio per
la classificazione nella maggior parte dei casi riesce a
distinguere correttamente router, PC, stampanti, telecamere e
dispositivi mobili.

Il modulo di trasferimento file aggiunge una funzionalità pratica
che supera il semplice ``vedere chi c'è in rete'' e lo rende
utile anche per operazioni quotidiane.

Essendo scritto interamente in Python, il programma è facile da
modificare e da estendere.  Il prezzo da pagare è che non sarà
mai veloce quanto un software scritto in C, ma per reti fino a
un paio di centinaia di dispositivi le prestazioni sono più che
sufficienti.

\end{document}
